\subsection{A bounded implementation pilot}\label{sec:pilot}

We measured construction and repeated-query costs in a single-threaded
CPython 3.15.1 implementation on an Apple M5 Pro with 48 GiB of memory,
running macOS 26.6.2 on AC power. Every method uses exact integers in
subnormal-ulp units. Thus a case with $(p,E)=(24,254)$ models the
numerical values of binary32 but does not execute native binary32
addition. The baseline replays the operand list and exits early on
absorbing overflow. Its input is already known to be representable.
The indices retain their representability check. We compare the original
sequential index constructor with balanced construction, and compare
replay queries with a query cache storing numerical piece origins.
The cache uses one binary search, shifts, and masks; it stores no operands.
Its preparation time and retained memory are charged separately.

The pilot has ten synthetic cases: small signed steps, paired cancellation,
and random exponents. Each uses 256 queries, half broadly distributed by
format rank and half near the operand scale. We retain five construction
samples and seven query samples after warmup, with randomized method order.
Each timed query batch repeats the fixed query set; batch calibration
targets at least 20 ms subject to a repetition cap. Profiling and allocation
tracing run separately. CPU frequency and affinity are not fixed, and
desktop background processes remain active. The saved results include
every timing sample, input, seed, and observed range.

\begin{table}[htbp]
\centering\small
\setlength{\tabcolsep}{4pt}
\begin{tabular}{llrrrrr}
\toprule
$(p,E,m)$ & Case & $T_s$ & $T_b$ & $t_r$ & $t_c$ & $Q_\ast$\\
\midrule
$(8,8,128)$ & Random & 10.85 & 27.21 & 1.681 & 0.215 & 18,563\\
$(24,254,1)$ & Steps & 12.33 & 2.06 & 0.409 & 0.412 & ---\\
$(24,254,32)$ & Steps & 436.69 & 177.55 & 9.530 & 0.412 & 19,516\\
$(24,254,128)$ & Cancel & 1844.52 & 696.70 & 36.883 & 0.424 & 19,119\\
$(53,2046,8)$ & Cancel & 3352.72 & 521.27 & 4.306 & 0.651 & 143,883\\
\bottomrule
\end{tabular}
\caption{Selected pilot medians. Sequential and balanced construction
times $T_s,T_b$ are in milliseconds; replay and cached-query times
$t_r,t_c$ are in microseconds. $Q_\ast$ includes cache preparation and is
an estimated crossover, not a measured end-to-end run. All queries in
the small random case overflow, and replay retains its early exit.}
\label{tab:pilot}
\end{table}

For a positive measured per-query saving, the linear cost model gives
\[
 Q_\ast=\left\lceil\frac{T_b+T_{\rm cache}}{t_r-t_c}\right\rceil,
\]
with times converted to the same units. No finite crossover exists in
this model when $t_r\leq t_c$. Table~\ref{tab:pilot} shows why a faster
query alone is insufficient: setup requires tens or hundreds of thousands
of queries to amortize in several broad-format cases. The one-addition
binary32-parameter case does not amortize. The corresponding binary64
case has only a roughly two-percent median saving, with overlapping
sample ranges; we do not infer a reliable advantage from it. The small
random case also shows that balanced construction can be slower than
the sequential constructor. These synthetic workloads do not establish
an application-level benefit or a speedup over native floating-point code.

Allocation measurements exposed a separate implementation cost: the
initial balanced constructor materializes all leaves. A lazy stack
instead merges adjacent equal-length blocks as they arrive and folds
the remaining blocks from right to left. It produces the same ordered
merge tree and final descriptors. In a separate paired run, traced peak
Python allocations for the $(24,254,128)$ cancellation case fell from
57.72 MB to 3.53 MB, and for $(53,2046,8)$ from 55.18 MB to 26.96 MB.
Construction times changed only within the observed variation. These
are reductions in temporary Python allocations, not process RSS or the
bit-encoding bound. The sequential constructor still used less peak
memory: about 0.69 MB and 9.66 MB in the respective cases. The query cache
also increases retained memory for the broad formats.

Before timing, all saved queries were checked against exact replay.
A separate enumerated-neighbour oracle checked 28,260 comparisons on
100 small-format blocks. Lazy construction passed another 7,092 oracle
comparisons on 108 blocks, including empty blocks, non-power-of-two
lengths, and one-pass operand iterators. These checks and identical
descriptors support the implementation changes; the scheduling and
caching optimizations are not additional mathematical novelty.
